Typically, ground fills may be helpful on lower-layer boards (1 or 2 layers) where good ground return paths are not available (slide 10 \cite{Kay2022Crosstalk}).
Since the processor board had ground layers adjacent to every signal layer, and no signal layers are adjacent to each other, it was deemed not necessary to include ground filling on the processor board \cite{Kay2022Crosstalk}. If ground filling was included, it would require several calculations to determine the required stitching via spacing, as well as a reconstruction of the signal layout to accomodate pours and fills inbetween each signal - a resource-costly step for minimal improvements.
\nl
Additionally, it was decided that the sensor board did not require ground filling. As there were several parallel routes in such close quarters, stitching vias could not be added, so any ground filling that wasn't correctly stitched would couple parallel lines. If the minimum neck length of the pour was increased to avoid long single-ended regions, then the ground fill would only pour on areas outside of the BGA, so including it would not have much effect.